\documentclass{article}
\usepackage[utf8]{inputenc}
\usepackage{hyperref}


%Remove below between comments to remove DRAFT watermark
% \usepackage[printwatermark]{xwatermark}
\usepackage{xcolor}
\usepackage{graphicx}
\usepackage{lipsum}

% \newwatermark[allpages,color=red!50,angle=45,scale=3,xpos=0,ypos=0]{DRAFT}
%Remove above this comment to remove DRAFT watermark

\title{Dash Simulator}
\author{Nakul Chawla}
% \date{February 2018}
\begin{document}

\maketitle

% \begin{abstract}
% Not long after the introduction of Nakamoto networks such as Bitcoin \cite{Nak}, there has been much debate about the possibility and ability of scaling these networks.  We estimate the throughput of these networks under a variety of different configurations and implementations.  Of particular emphasis, we study how different block propagation techniques effect this throughput limit.
% \end{abstract}

\section{Introduction}
Our Bitcoin simulator is an NS3 based
simulator that was originally built for the research shown in
\cite{security}. The simulator demonstrates the workings of Bitcoin, Dogecoin and
Litecoin, but bases its research on each block being the fundamental/rudimentary component. 

Below are the modifications that have been done to the bitcoin simulator in order for us to experiment newer propagation
technologies.

\section{Modifications}
We found that simulations that merely went down to the block level were not detailed enough for our purpose.  In order to reliably model a consensus network with block propagation speedups such as compact or xthin we must understand these networks down at the level of transactions.

In order to understand the network at the level of transactions, we modeled generation of transactions using a normal
distribution. Each node is has a transaction generator that creates a random transaction with a piecewise constant
distribution and hashes that data.

There are 3 types of nodes : Full nodes, Miners and Masternodes. Masternode in the current simulator is
not a separate class but a check on full node.
The network is randomly distributed with full nodes, masternodes and the miners. All the miners in the network are
connected to each other which simulates the real world network.
When establishing the connections between nodes, we create point-to-point channels between them, which abstracts away
any intermediate devices (routes, switches, etc). These channels have two characteristics: 
\begin{enumerate}
    \item latency
    \item bandwidth
\end{enumerate}

Mining pools typically maintain private peering connections among
themselves—which we capture in our simulations. Besides direct peering, a
number of mining pools nowadays participate in Matt Corallo’s relay network
\cite{matt} that is operated independently of the default Bitcoin P2P overlay network.


In short, below are the modifications done to the bitcoin-simulator for it to run the current dash-simulations:

\begin{enumerate}
    \item Transactions are added in order to add traffic and understand propagations in more detail.
    \item Every node now has it's own mempool in order to test set reconciliation for xthin propagation.
    \item Compact block propagation is added.
    \item Extreme thin block propagation is added.
    \item Masternode (bool) is added to the full-node. Because this is network propation, the bandwidth is the only
        differentiating factor between a full node and masternode. 
    \item Bloom filter using Murmur Hash is added in order to test xthin.
    \item Capability to test x-11 hashing is added, although it doesn't add value to propagation. 
    \item Fill Block - This parameter if true, fills the block with transactions to its maximum capacity. If this is set
        to false, the simulator uses a piecewise constant distribution to generate the next block size and adds
        transactions limiting to the new block size. 
\end{enumerate}

The resulting simulator code is available at \cite{sim}.  

\section{Parameters}
The simulator takes in the below given parameters:

\begin{enumerate}
    \item Block Size
    \item Block interval
    \item Propagation type
    \item Number of blocks
    \item Fill Block  - The block is filled to capacity if fillBlock=true.
\end{enumerate}

\section{Results and Findings}

We ran multiple experiments in order to find how to scale the dash network for future:

\subsection{Orphan rate based on type of block}
We tested 3 types of blocks :

\begin{enumerate}
    \item Traditional
    \item Compact : BIP 152
    \item Extreme Thin Blocks (xthin) : BUIP 10
\end{enumerate}

The simulations were run with different network sizes and different block sizes.
Through these simulations, it was realized that both compact and xthin blocks can support the Dash network with a
capacity an order of magnitude larger than the traditional block propagation.

However, the important decision is to compare the propagation between Compact and xthin blocks and which will work
better for the future of Dash network.

The first set of experiments was run on a network of size 24 with 16 mining pools and 8 MasterNodes.
For this small network, compact blocks at 160 MB resulted in an orphan rate in the ballpark of 20\% where as xthin
blocks produced orphan rates less than 5\% for block size 320MB. 

In order to test the current network size, we ran experiments with 6000 nodes comprising of 8 mining pools and 5000
masternodes. Compact blocks produced consistently high orphan rates where as xthin ranged between 3\% and 15\%. 

Since, there is high variance in orphan rates for xthin, with some simulations resulting in lower orphan rates as
compared to compact blocks which consistently produced higher orphan rates, it is currently the best propagation method for dash to adopt.

The results of the simulator are produced in the form of log-files that can be found at \cite{results}.

\subsection{Role of Masternodes}
The current dash network is a 2-tier network with Masternodes that maintain higher hardware capabilities for better
propagation and, faster transaction verification and validation.

We know that Masternodes add two features to the dash network:

\begin{enumerate}
    \item Privatesend
    \item Instantsend
\end{enumerate}

Both of these are not available on any other network except for dash. But what we realized through our simulations that
for block sizes in the ballpark of 10MB, Masternodes don't add significant value to the network in reducing orphan rate.

The two major reasons to that is :

\begin{enumerate}
    \item The hardware specification of a full-node is capable of verifying and validating transactions at a fast
        pace for block sizes in the neighbourhood of 10MB , orphan rate is not affected. 
    \item For bandwidth, we ran an experiment with all masternodes running at 100MB per second bandwidth, and the orphan
        rate did not change significantly.
\end{enumerate}

Larger block sizes close to 500MB may be affected by masternodes, especially in propagation if the network requirements
are maintained at 1GB per second bandwidth with higher up-time.  

  \bibliographystyle{alpha}
  \bibliography{bib.bib}
\end{document}
